\section{Introduction and current theorem structure}\label{sec:current-overview}\label{sec:introduction}
One physical command and several commands pose different realization problems. A stochastic transition matrix may reproduce one complete output sequence while no common finite family can reproduce every commanded sequence on the same retained labels. Likewise a periodic output's period need not equal its smallest stochastic memory. The present revision addresses both questions while keeping a prescribed initial label and a bounded terminal decoder.

Theorem~\ref{thm:r72-reversible} characterizes exact finite realization for every compactly reversible finite-dimensional response: the observable response group must be finite, and that condition is sufficient. The necessity proof first intersects the hidden simplex with the reachable linear space, then projects that polytope onto the response space. This order is essential. Compact reversibility makes its inclusions into equalities, so a faithful finite vertex permutation forces a finite group. Corollary~\ref{cor:r72-unitary} applies the criterion to prepared unitary command families without requiring that the protected physical state be observed.

Theorem~\ref{thm:r72-control-envelope} turns this qualitative obstruction into a matching quantitative law for two noncommuting involutions. One mixed-word subsequence provides a necessary stochastic eigenvalue, while one shared polygon supplies both command matrices. The resulting arithmetic envelope has the same two exact growth exponents and cubic-root criterion as the inherited single-frequency theory. Testing each command separately would miss this obstruction. The theorem concerns exact response to arbitrary supplied commands, not an unrestricted optimal-feedback classification.

Theorem~\ref{thm:r72-periodic-memory} gives a different exact conclusion for a finite rational-frequency mixture. It minimizes a sum of least-common-multiple cycle sizes over partitions of the actual spectral support, adding a counted initialization label when several recurrent classes are needed. Positive coefficients permit a bounded decoder on each constructed cycle; arbitrary real stochastic competitors satisfy the same lower through their closed classes. Orders $3,5,7$ require exactly $16$ stochastic labels but $105$ deterministic-transition labels. Theorem~\ref{thm:r72-periodic-leakage} preserves the complete partition lower through a decaying output-visible core and proves the exact additional live label for physical leakage. Corollary~\ref{cor:r72-synthesis} supplies finite minimal synthesis and a constant-memory zero-noise plateau.

These arguments build on established invariant-cone realization and stochastic spectral theory \cite{Heller65,MonrasWinter14,ThonJaeger15,Nussbaum98,BenvenutiFarina04}. Finite-dimensional quantum processes without finite classical models, as well as noise-dependent classical-memory lower bounds, predate this article \cite{Fanizza24}. We do not claim those general phenomena or the classical cycle method as new. The present assertions are the compact observable-group criterion, the common all-word arithmetic realization, and the exact bounded-response partition, initialization and leakage laws. Their independent novelty and significance are not certified by a literature search.

The complete predecessor Diophantine envelope and exponent proofs remain below. Their continued-fraction and stochastic-sector ingredients retain their credit and explicit proof \cite{DmitrievDynkin46,JohnsonPaparella17,KirklandSmigoc22,HoneVarona19}. The actual-history causal transfer also remains in the main article. Its accepted comparison measure is formed on the competing experiment's actual pre-call histories, not by assuming that the comparison experiment supports a terminating reference run.

The current supplement retains the full squared-loss embedding, normalized-budget stopping rigidity, robust tied optimality, state-interior instrument realization, performed-law duality, reset, rank, timing and geometric theory. In the unary squared-task interpretation no protected-state oracle, free initial mixture or uncharged persistent seed is introduced. Controlled simulation labels, observer labels, terminal values, program storage and protected physical dimension remain distinct. The original research identity is unchanged; the new theorems neither silently delete independent results nor turn them into unproved corollaries.
